\documentclass[10pt,a4paper]{amsart}
\usepackage[margin=19mm]{geometry}
\usepackage{amsmath,amssymb,mathtools}
\usepackage[scr=rsfs,cal=euler]{mathalfa}
\usepackage{xfrac}
\usepackage[unicode,hidelinks,bookmarks=false]{hyperref}
\newcommand{\ie}{i.e.\ }
\newcommand{\cf}{cf.\ }
\newcommand{\resp}{respectively\ }
\newcommand{\Subsection}{\subsection}
\providecommand{\nobreakdash}{\nobreak\hbox{-}}
\setlength{\emergencystretch}{2em}
\multlinegap0pt
\usepackage{dsfont}
\usepackage{amssymb}
\usepackage{mathtools}
\usepackage{caption}
\usepackage[shortlabels]{enumitem}
\usepackage{tikz}
\usepackage[all]{xy}
\usepackage{xcolor}
\newtheorem{proposition}{Proposition}
\DeclareMathOperator{\Frac}{Frac}
\newcommand{\HH}{\mathrm{H}}
\DeclareMathOperator{\Reg}{Reg}
\newcommand{\ff}{\mathrm{f}}
\title{On generalized main conjectures and $p$-adic Stark conjectures for Artin motives}
\author{Alexandre Maksoud}
\date{Source: arXiv:2103.06864v4 (CC BY 4.0)}
\begin{document}
\maketitle

\noindent\emph{Source and licence.} On generalized main conjectures and $p$-adic Stark conjectures for Artin motives. Version arXiv:2103.06864v4 (CC BY 4.0), distributed by arXiv under the Creative Commons Attribution 4.0 International licence, http://creativecommons.org/licenses/by/4.0/, as recorded in the arXiv licence metadata for this exact version and shown on the abstract page https://arxiv.org/abs/2103.06864v4. Full licence notice: \texttt{licenses/arxiv-2103.06864-CC-BY-4.0.txt}.\par\smallskip

\noindent\textit{Editorial scope.} Counted unit: the article's proposition prop:changement\_de\_base (source0.tex lines 1311--1313). The article's complete original proof of it is delivered with it (source0.tex lines 1314--1324, one \texttt{proof} environment of 11 lines). No mathematics is written, repaired or re-derived: the statement and the proof are the article's own lines, in the article's own notation, and no retained line is substituted or re-flowed.  No further source-local context is needed: every cross-reference of every retained line resolves inside the two delivered spans.  Nothing was omitted from any retained span, so the package carries no omission record.  No retained line contains a citation, so the wrapper carries no bibliography and no bibliography entry.  No retained line contains a figure environment, an image-inclusion command or drawing code, so no image is delivered and the package declares no asset dependency.  Editorial formatting only, in addition: an AMS \texttt{amsart} preamble that restates the article's own declarations and macros used by the retained lines, namely the environments \texttt{proposition} on separate counters, together with the standard packages those lines need.  The retained environments print in this document as Proposition 1 in order of appearance, whereas the article prints the same items with the numbering of its own full document.  The complete native source of the article as distributed in the arXiv e-print for this exact version is bundled unchanged as \texttt{sources/109065/source0.tex}.
\begin{proposition}\label{prop:changement_de_base}
	If \textbf{EX}$_{\rho,\rho^+_\alpha}$ holds for all $\alpha\in I_{\rho^+}$, then \textbf{EX}$_{\rho,\rho^+}$ holds as well, and $\theta_{\rho,\rho^+} = \sum_{\alpha \in I_{\rho^+}}c_\alpha\cdot \theta_{\rho,\rho^+_\alpha}$.
\end{proposition}

\begin{proof}
	Assume that \textbf{EX}$_{\rho,\rho^+_\alpha}$ holds for all $\alpha\in I_{\rho^+}$. For every character $\eta\in\widehat{\Gamma}$, the rule $\omega_p^+ \mapsto \Reg_{\omega_p^+}(\rho\otimes\eta)$ (where the $p$-adic regulator is computed in a fixed basis $\omega_{\ff,\eta}$ of $\det_E \HH^1_\ff((\rho\otimes\eta)^*(1))$) defines a $E_{p}$-linear map 
	${\bigwedge}^{d^+} W_p \longrightarrow E_{p,\eta},$ so we have $$\Reg_{\omega_p^+}(\rho\otimes\eta)=\sum_{\alpha \in I_{\rho^+}} c_\alpha\cdot \Reg_{\omega^+_{p,\alpha}}(\rho\otimes\eta).$$ Therefore, the element $\theta_{\rho,\rho^+}\in\Frac(\Lambda)$ defined as $\sum_{\alpha \in I_{\rho^+}}c_\alpha\cdot \theta_{\rho,\rho^+_\alpha}$ satisfies
%	\begin{align*}
%	\eta(\theta_{\rho,\rho^+}) &= \sum_{\alpha \in I_{\rho^+}}c_\alpha\cdot \eta(\theta_{\rho,\rho^+_\alpha}) \\
%	&=M_{\rho,\eta} \cdot \sum_{\alpha \in I_{\rho^+}}c_\alpha\cdot \frac{\Reg_{\omega^+_{p,\alpha}}(\rho\otimes\eta)}{\det(\rho_\alpha^-)(\sigma_p^n))} \\
%	&=M_{\rho,\eta} \cdot \frac{\Reg_{\omega^+_{p}}(\rho\otimes\eta)}{\det(\rho^-)(\sigma_p^n))}
%	\end{align*}
	\[\eta(\theta_{\rho,\rho^+}) = \sum_{\alpha \in I_{\rho^+}}c_\alpha\cdot \eta(\theta_{\rho,\rho^+_\alpha}) =M_{\rho,\eta} \cdot \sum_{\alpha \in I_{\rho^+}}c_\alpha\cdot \frac{\Reg_{\omega^+_{p,\alpha}}(\rho\otimes\eta)}{\det(\rho_\alpha^-)(\sigma_p^n))} =M_{\rho,\eta} \cdot \frac{\Reg_{\omega^+_{p}}(\rho\otimes\eta)}{\det(\rho^-)(\sigma_p^n))}\]
	for all non-trivial characters $\eta\in\widehat{\Gamma}$ of conductor $p^n$, where we have put $M_{\rho,\eta}=\frac{\tau(\eta)^{d^-}}{\tau(\rho\otimes\eta)}\ \frac{L^{*}\left((\rho \otimes \eta)^*,0\right)}{\Reg_{\omega_\infty^+}(\rho\otimes\eta)}$. Therefore, $\theta_{\rho,\rho^+}$ satisfies the interpolation property of \textbf{EX}$_{\rho,\rho^+}$. Since it has no pole except maybe at $\mathds{1}$ by construction, we have shown that \textbf{EX}$_{\rho,\rho^+}$ is valid.
\end{proof}

\end{document}
